% Ported from an earlier report section (one-time conversion; edit here).
\subsection{More explorations of the campaign data \texorpdfstring{\screen}{[screen]}}\label{vt:sec:more}

Everything in this section is exploratory: it uses all \NScenarios\ scenarios, was not preregistered, and is
meant to show \emph{how} the costs arise. It uses \PerTurnNFmt\ turn rows: the \NTurns\ turn rows of the campaign minus the \RvKilledTurns\ turns of the
memory-killed session and \RvNSkips\ final turns that were skipped in sessions that otherwise stayed cost-valid
(\RvSkipList). Each of those two sessions was billed for one turn fewer than its partner, a negligible pair bias.

\subsubsection{How cost accumulates over a session}

\begin{figure}[tbp]
\centering
\pgfplotslegendfromname{perturnlegend}\par\smallskip
% Mean cost per turn by turn index, per arm, for each host (all splits, cost-valid sessions).
\begin{tikzpicture}
\begin{groupplot}[benchmark, group style={group size=2 by 1, horizontal sep=1.6cm},
  width=0.48\linewidth, height=5.6cm, xmin=0.5, xmax=16.5, xtick={1,4,8,12,16}, ymin=0, ymajorgrids,
  xlabel={Turn}, title style={font=\small\bfseries},
  cycle list={{okgray, mark=*, solid}, {okorange, mark=square*, dashed}, {okblue, mark=triangle*, solid}, {okgreen, mark=diamond*, densely dotted}},
  every axis plot/.append style={line width=1pt, mark size=1.8pt, unbounded coords=jump}]
\nextgroupplot[title={Fable host}, ylabel={Mean cost of the turn (\$)},
  legend to name=perturnlegend, legend columns=4, legend cell align=left,
  legend style={/tikz/every even column/.append style={column sep=8pt}}]
\addplot+ table[x=turn, y=anchor]{generated/data/perturn-fable.dat}; \addlegendentry{plain host}
\addplot+ table[x=turn, y=shipped]{generated/data/perturn-fable.dat}; \addlegendentry{shipped (re-chooses each turn)}
\addplot+ table[x=turn, y=sticky]{generated/data/perturn-fable.dat}; \addlegendentry{sticky}
\addplot+ table[x=turn, y=sonnet]{generated/data/perturn-fable.dat}; \addlegendentry{Sonnet only}
\nextgroupplot[title={Opus host}]
\addplot+ table[x=turn, y=anchor]{generated/data/perturn-opus.dat};
\addplot+ table[x=turn, y=shipped]{generated/data/perturn-opus.dat};
\addplot+ table[x=turn, y=sticky]{generated/data/perturn-opus.dat};
\addplot+ table[x=turn, y=sonnet]{generated/data/perturn-opus.dat};
\end{groupplot}
\end{tikzpicture}

\caption{On Fable the routed arms are cheaper at every turn; on Opus they cost more at every turn. Mean cost of each turn by turn index, per arm and host. Later turns are averaged over fewer scenarios,
because scenarios have \MinTurns\ to \MaxTurns\ turns.}
\label{vt:fig:perturn}
\howtoread{Each line is one arm. Turn 1 is expensive for everyone (a cold cache); later turns are cheaper but
grow as the history grows. On Fable the routed lines stay well below the plain host at every turn; on Opus they
sit above it.}
\end{figure}

On the Fable host the first turn alone was \TurnOneShareFable\,\% of a plain session's cost on average.
\Cref{vt:fig:turncomp} shows what each turn of a plain-host session paid for.

\begin{figure}[tbp]
\centering
\pgfplotslegendfromname{turncomplegend}\par\smallskip
% What each turn of a plain-host session paid for, by turn index (mean dollars per class).
\begin{tikzpicture}
\begin{groupplot}[benchmark, group style={group size=2 by 1, horizontal sep=1.6cm},
  width=0.48\linewidth, height=5.6cm, xmin=0.3, xmax=16.7, xtick={1,4,8,12,16},
  ymin=0, ymajorgrids, xlabel={Turn}, title style={font=\small\bfseries}]
\nextgroupplot[ybar stacked, bar width=5pt, title={Fable host, plain}, ylabel={Mean cost of the turn (\$)},
  legend to name=turncomplegend, legend columns=4, legend cell align=left,
  legend style={/tikz/every even column/.append style={column sep=8pt}}]
\addplot[fill=okgray!50, draw=black!60] table[x=turn, y=input]{generated/data/turncomp-fable.dat}; \addlegendentry{uncached input}
\addplot[fill=okblue!45, draw=black!60] table[x=turn, y=read]{generated/data/turncomp-fable.dat}; \addlegendentry{cache read}
\addplot[fill=okorange!80, draw=black!60, postaction={pattern=north east lines, pattern color=white}] table[x=turn, y=write]{generated/data/turncomp-fable.dat}; \addlegendentry{cache write}
\addplot[fill=okgreen!70, draw=black!60, postaction={pattern=dots, pattern color=white}] table[x=turn, y=output]{generated/data/turncomp-fable.dat}; \addlegendentry{output}
\nextgroupplot[ybar stacked, bar width=5pt, title={Opus host, plain}]
\addplot[fill=okgray!50, draw=black!60] table[x=turn, y=input]{generated/data/turncomp-opus.dat};
\addplot[fill=okblue!45, draw=black!60] table[x=turn, y=read]{generated/data/turncomp-opus.dat};
\addplot[fill=okorange!80, draw=black!60, postaction={pattern=north east lines, pattern color=white}] table[x=turn, y=write]{generated/data/turncomp-opus.dat};
\addplot[fill=okgreen!70, draw=black!60, postaction={pattern=dots, pattern color=white}] table[x=turn, y=output]{generated/data/turncomp-opus.dat};
\end{groupplot}
\end{tikzpicture}

\caption{Turn 1 is mostly cache writes; later turns are mostly cache reads. What each turn of a plain-host session paid for, by turn index (mean dollars per token class).}
\label{vt:fig:turncomp}
\howtoread{Turn 1 is mostly cache writes (hatched): the whole prompt is new. Later turns are dominated by cache
reads (blue) as the history is re-read each time, and the writes shrink to the new part of the conversation.
Fable's writes and output are expensive; Opus's reads are cheap.}
\end{figure}

\subsubsection{Long pauses and the cache}

A pause longer than the cache lifetime should force the next request to write the whole history again.
\Cref{vt:fig:gapwrites} confirms it: after a \LongGapMin-minute pause, a plain Fable session's first request wrote
\GapWriteFableLong\ tokens on average, against \GapWriteFableShort\ after a \ShortGapS-second pause
(\GapWriteFableX\X; \GapWriteOpusX\X\ on Opus).

\begin{figure}[tbp]
\centering
\pgfplotslegendfromname{gaplegend}\par\smallskip
% Cache writes on the first request of a turn, after a short pause vs after a long (expiring) pause.
\begin{tikzpicture}
\begin{axis}[benchmark, xbar, bar width=6pt, scale only axis, width=10.5cm, height=5cm, xmin=0,
  ymin=-0.6, ymax=4.6, ytick={0,...,4}, yticklabels from table={generated/data/gapwrites.dat}{label},
  xlabel={Tokens written on the turn's first request (thousands)}, xmajorgrids,
  title={A pause longer than the cache lifetime}, title style={font=\small\bfseries},
  legend to name=gaplegend, legend columns=2, legend cell align=left,
  legend style={/tikz/every even column/.append style={column sep=10pt}}]
\addplot[fill=okblue!45, draw=black!60] table[x=short, y=idx]{generated/data/gapwrites.dat};
\addlegendentry{after a \ShortGapS-second pause}
\addplot[fill=okorange!80, draw=black!60, postaction={pattern=north east lines, pattern color=white}] table[x=long, y=idx]{generated/data/gapwrites.dat};
\addlegendentry{after a \LongGapMin-minute pause}
\end{axis}
\end{tikzpicture}

\caption{After a long pause the first request rewrites the history on every arm and host. Cache-write tokens on the first request of a turn, after a short pause and after a long pause (turns 2
onwards), per arm and host.}
\label{vt:fig:gapwrites}
\howtoread{Each row has two bars. The hatched bar (after a long pause) is several times longer in every row:
the expired history is written again, on any model.}
\end{figure}

\subsubsection{When the shipped arm switched}

The shipped arm re-chooses its model at the start of each turn. Its \SwitchTotal\ switches were all at turn
boundaries; \SwitchTurnTwoPct\,\% of them led into turn 2 (\cref{vt:fig:switch}). Each switch cost about
\RebuildPerSwitchCents\ cents of rebuild on average.

\begin{figure}[tbp]
\centering
\pgfplotslegendfromname{switchlegend}\par\smallskip
% When the shipped arm switched models: switches into each turn, by host.
\begin{tikzpicture}
\begin{axis}[benchmark, ybar, bar width=5pt, scale only axis, width=11cm, height=4.4cm, xmin=1.3, xmax=16.7,
  xtick={2,...,16}, ymin=0, xlabel={Turn the switch led into}, ylabel={Switches}, ymajorgrids,
  title={Shipped arm: model switches by turn}, title style={font=\small\bfseries},
  legend to name=switchlegend, legend columns=2, legend cell align=left,
  legend style={/tikz/every even column/.append style={column sep=10pt}}]
\addplot[fill=okblue!60, draw=black!60] table[x=turn, y=fable]{generated/data/switch-timing.dat};
\addlegendentry{Fable host}
\addplot[fill=okorange!80, draw=black!60, postaction={pattern=north east lines, pattern color=white}] table[x=turn, y=opus]{generated/data/switch-timing.dat};
\addlegendentry{Opus host}
\end{axis}
\end{tikzpicture}

\caption{The shipped router's model switches are spread over the whole session. Shipped arm: number of model switches leading into each turn, by host.}
\label{vt:fig:switch}
\howtoread{Each pair of bars is one turn. Switches are spread over the whole session rather than concentrated
at the start; every one of them re-writes the history to the other model's cache.}
\end{figure}


\subsubsection{Scenario by scenario}

Averaging hides how uneven the saving is. Most scenarios saved money on the Fable host, and the expensive ones
tended to save the most; \ScNegative\ of \ScN\ scenarios cost more with sticky routing on average.
\Cref{vt:fig:scenario} shows each scenario once.

\begin{figure}[tbp]
\centering
\pgfplotslegendfromname{scenariolegend}\par\smallskip
% Scenario-level view (Fable host, sticky): mean plain-host cost vs mean saving per session.
\begin{tikzpicture}
\begin{axis}[benchmark, scale only axis, width=\ScAxisW, height=\ScAxisH, xmin=0, xmax=\ScXMax, xtick distance=2, ytick distance=2,
  ymin=\ScYMin, ymax=\ScYMax, xlabel={Plain Fable session cost (\$, mean of 2 repetitions)},
  ylabel={Saving of sticky (\$ per session)}, xmajorgrids, ymajorgrids, clip=false,
  title={One point per scenario}, title style={font=\small\bfseries},
  legend to name=scenariolegend, legend columns=2, legend cell align=left,
  legend style={/tikz/every even column/.append style={column sep=10pt}}]
\addplot[okgray, dashed, line width=0.8pt, domain=0:\ScXMax, samples=2] {0};
\addlegendentry{no saving}
\addplot[only marks, mark=*, mark size=2pt, color=okblue] table[x=anchor, y=saving]{generated/data/scenario-scatter.dat};
\addlegendentry{scenario (sticky arm, Fable host)}
% Generated by build_assets.py: labels for panel scenario-scatter. Do not edit.
\node[scatterlabel, anchor=east] at ($(axis cs:7.7791,-4.1312)+(-4.700pt,0.000pt)$) {arch-explain};
\node[scatterlabel, anchor=west] at ($(axis cs:9.2815,-3.2406)+(4.700pt,0.000pt)$) {py-go-counting};
\node[scatterlabel, anchor=south] at ($(axis cs:12.3685,-2.6644)+(0.000pt,4.700pt)$) {py-sgf-parsing};

\end{axis}
\end{tikzpicture}

\caption{Most scenarios saved money on Fable, and expensive scenarios tended to save more. Scenario by scenario: the mean plain Fable cost against the mean saving of the sticky arm (Fable host). The three
scenarios where sticky cost the most extra are labelled.}
\label{vt:fig:scenario}
\howtoread{Points above the dashed line saved money. Most scenarios did, and expensive scenarios tended to save
more. \ScNegative\ of \ScN\ scenarios cost more on average; the labelled ones are the worst three.}
\end{figure}


\subsubsection{The noise floor in detail}

The A/A pairs show how far two identical sessions drift apart. \Cref{vt:fig:aahist} shows the whole distribution
behind the summary in \cref{vt:sec:aa}: most pairs fall within a few per cent, a handful differ by up to a quarter, and
the Opus distribution sits slightly to the right of zero, a bias we cannot explain (\cref{vt:sec:aa}).

\begin{figure}[tbp]
\centering
\pgfplotslegendfromname{aalegend}\par\smallskip
% A/A: distribution of log cost ratios between two identical plain-host sessions.
\begin{tikzpicture}
\begin{axis}[benchmark, ybar, bar width=6pt, scale only axis, width=10cm, height=4.4cm, xmin=-0.3, xmax=0.3,
  ymin=0, xlabel={Log cost ratio, second plain session over first}, ylabel={Pairs}, ymajorgrids,
  title={Two identical sessions}, title style={font=\small\bfseries},
  legend to name=aalegend, legend columns=2, legend cell align=left,
  legend style={/tikz/every even column/.append style={column sep=10pt}}]
\addplot[fill=okblue!60, draw=black!60] table[x=mid, y=fable]{generated/data/aa-hist.dat};
\addlegendentry{Fable host}
\addplot[fill=okorange!80, draw=black!60, postaction={pattern=north east lines, pattern color=white}] table[x=mid, y=opus]{generated/data/aa-hist.dat};
\addlegendentry{Opus host}
\draw[black!70, line width=0.8pt] (axis cs:0,0) -- (axis cs:0,\pgfkeysvalueof{/pgfplots/ymax});
\end{axis}
\end{tikzpicture}

\caption{Two identical sessions differ by chance; this spread is the noise floor. Distribution of the log cost ratio between two identical plain-host sessions (A/A), by host.}
\label{vt:fig:aahist}
\howtoread{A perfectly repeatable system would put every bar at zero. The spread is the noise floor; the Opus
bars lean slightly to the right, the small bias discussed in \cref{vt:sec:aa}.}
\end{figure}


\subsubsection{Quality by scenario family, and savings next to quality}\label{vt:sec:joint}

\Cref{vt:tab:qualityfamily} gives the turn-pass difference by family; \cref{vt:tab:jointfamily,vt:tab:jointtask} put each
cell's saving per 1,000 sessions (as in \cref{vt:tab:per1000}) next to its turn-pass difference. The family with the
largest Fable sticky saving, knowledge, also has the largest quality loss: \RvKnowFableSticky\ (95\,\% CI
\RvKnowFableStickyCI), below the \NIMargin\ margin. Inside knowledge the losses come from the explain and review
scenarios. This is a family effect, not a general rule that bigger savings cost quality: across scenarios, the rank
correlation between saving and turn-pass difference for Fable sticky is \RvSpearman.

\begin{table}[tbp]
\centering\small
\caption{Mean difference in turn-pass fraction (arm minus plain host) by scenario family, all splits. Negative
means the arm passed fewer scripted turns.}
\label{vt:tab:qualityfamily}
\begin{tabular*}{\textwidth}{@{\extracolsep{\fill}}l r r r r r r r@{}}
\toprule
 & & \multicolumn{3}{c}{Fable host} & \multicolumn{3}{c}{Opus host} \\
\cmidrule(lr){3-5}\cmidrule(l){6-8}
Family & scenarios & sticky & shipped & sonnet & sticky & shipped & sonnet \\
\midrule
knowledge & 14 & $-$0.083 & $-$0.061 & $-$0.029 & $-$0.060 & $-$0.048 & $-$0.018 \\
mixed & 20 & 0.009 & 0.001 & 0.015 & 0.010 & $-$0.005 & 0.019 \\
polyglot & 20 & $-$0.044 & $-$0.025 & $-$0.041 & $-$0.007 & $-$0.020 & 0.000 \\
repos & 16 & 0.025 & 0.026 & 0.023 & 0.005 & 0.017 & 0.016 \\
\bottomrule
\end{tabular*}

\end{table}

\begin{table}[tbp]
\centering\scriptsize
\caption{Savings and quality together, by family: saving per 1,000 sessions (US dollars) / mean turn-pass
difference, all splits. $\dagger$: the mean difference is below the \NIMargin\ margin.}
\label{vt:tab:jointfamily}
\begin{tabular*}{\textwidth}{@{\extracolsep{\fill}}l r r r r r r@{}}
\toprule
 & \multicolumn{3}{c}{Fable host} & \multicolumn{3}{c}{Opus host} \\
\cmidrule(lr){2-4}\cmidrule(l){5-7}
Family & sticky & shipped & sonnet & sticky & shipped & sonnet \\
\midrule
knowledge & 2{,}973 / $-$0.083\textsuperscript{$\dagger$} & 2{,}727 / $-$0.061\textsuperscript{$\dagger$} & 2{,}576 / $-$0.029 & 31 / $-$0.060\textsuperscript{$\dagger$} & $-$93 / $-$0.048 & $-$403 / $-$0.018 \\
mixed & 1{,}783 / 0.009 & 1{,}517 / 0.001 & 2{,}024 / 0.015 & $-$126 / 0.010 & $-$289 / $-$0.005 & $-$351 / 0.019 \\
polyglot & 2{,}229 / $-$0.044 & 2{,}106 / $-$0.025 & 2{,}897 / $-$0.041 & $-$751 / $-$0.007 & $-$740 / $-$0.020 & $-$1{,}696 / 0.000 \\
repos & 1{,}047 / 0.025 & 761 / 0.026 & 739 / 0.023 & $-$1{,}070 / 0.005 & $-$1{,}203 / 0.017 & $-$1{,}883 / 0.016 \\
\bottomrule
\end{tabular*}

\end{table}

\begin{table}[tbp]
\centering\scriptsize
\caption{Savings and quality together, by task type (same layout as \cref{vt:tab:jointfamily}). Docs has no test-split
scenario (\cref{vt:tab:tasksplit}).}
\label{vt:tab:jointtask}
\begin{tabular*}{\textwidth}{@{\extracolsep{\fill}}l r r r r r r@{}}
\toprule
 & \multicolumn{3}{c}{Fable host} & \multicolumn{3}{c}{Opus host} \\
\cmidrule(lr){2-4}\cmidrule(l){5-7}
Task type & sticky & shipped & sonnet & sticky & shipped & sonnet \\
\midrule
bugfix & 1{,}346 / 0.027 & 797 / 0.025 & 871 / 0.026 & $-$1{,}037 / 0.010 & $-$1{,}240 / 0.017 & $-$1{,}948 / 0.019 \\
feature & 2{,}183 / $-$0.071\textsuperscript{$\dagger$} & 2{,}019 / $-$0.033 & 2{,}895 / $-$0.056\textsuperscript{$\dagger$} & $-$574 / $-$0.031 & $-$490 / $-$0.044 & $-$1{,}099 / $-$0.019 \\
mixed & 2{,}020 / 0.006 & 1{,}941 / $-$0.001 & 2{,}405 / 0.001 & $-$303 / 0.009 & $-$401 / $-$0.002 & $-$815 / 0.015 \\
docs & 2{,}615 / $-$0.033 & 2{,}263 / $-$0.010 & 2{,}246 / 0.008 & 264 / 0.019 & 162 / 0.023 & $-$129 / 0.035 \\
explain & 2{,}315 / $-$0.143\textsuperscript{$\dagger$} & 2{,}244 / $-$0.161\textsuperscript{$\dagger$} & 2{,}367 / $-$0.017 & 341 / $-$0.091\textsuperscript{$\dagger$} & 235 / $-$0.091\textsuperscript{$\dagger$} & 242 / 0.017 \\
review & 3{,}007 / $-$0.155\textsuperscript{$\dagger$} & 3{,}011 / $-$0.093\textsuperscript{$\dagger$} & 2{,}389 / $-$0.081\textsuperscript{$\dagger$} & $-$425 / $-$0.138\textsuperscript{$\dagger$} & $-$642 / $-$0.106\textsuperscript{$\dagger$} & $-$1{,}193 / $-$0.081\textsuperscript{$\dagger$} \\
\bottomrule
\end{tabular*}

\end{table}

\begin{table}[tbp]
\centering\small
\caption{Fable sticky by task type, all splits, with 95\,\% scenario-cluster bootstrap intervals on the saving per
1,000 sessions (US dollars) and on the turn-pass difference. $\dagger$: below the \NIMargin\ margin.}
\label{vt:tab:taskfable}
\begin{tabular*}{\textwidth}{@{\extracolsep{\fill}}l r r r r r@{}}
\toprule
Task type & scenarios & saving / 1,000 & 95\,\% CI & turn-pass $\Delta$ & 95\,\% CI \\
\midrule
bugfix & 18 & 1{,}346 & 759 to 1{,}971 & 0.027 & $-$0.004 to 0.080 \\
feature & 13 & 2{,}183 & 893 to 3{,}421 & $-$0.071\textsuperscript{$\dagger$} & $-$0.225 to 0.054 \\
mixed & 28 & 2{,}020 & 1{,}075 to 2{,}832 & 0.006 & $-$0.019 to 0.030 \\
docs & 4 & 2{,}615 & 1{,}896 to 3{,}388 & $-$0.033 & $-$0.072 to 0.000 \\
explain & 3 & 2{,}315 & 1{,}912 to 2{,}804 & $-$0.143\textsuperscript{$\dagger$} & $-$0.222 to $-$0.056 \\
review & 4 & 3{,}007 & 2{,}332 to 3{,}585 & $-$0.155\textsuperscript{$\dagger$} & $-$0.263 to $-$0.037 \\
\bottomrule
\end{tabular*}

\end{table}


\subsubsection{Requests per session}

Sessions on Sonnet needed more main-loop requests for the same scripted work than plain-host sessions did, and
that volume is part of why routing cost more on the Opus host. The sticky arm made fewer requests than its mix of
Sonnet and host sessions would suggest (\cref{vt:sec:effort}). \Cref{vt:fig:requests} shows the spread.

\begin{figure}[tbp]
\centering
% Main-loop requests per session (box: quartiles, line: median, whiskers: 5th-95th percentile).
\begin{tikzpicture}
\begin{axis}[benchmark, scale only axis, width=10.5cm, height=5.6cm, xmin=0, xmajorgrids,
  ytick={1,...,\ReqBoxN}, yticklabels/.expanded={\ReqBoxLabels},
  xlabel={Requests per session}, title={How many requests a session made}, title style={font=\small\bfseries},
  boxplot/draw direction=x, every axis plot/.append style={line width=0.9pt}]
\addplot+[boxplot prepared={lower whisker=22.0, lower quartile=36.0, median=47.5, upper quartile=58.0, upper whisker=78.0}, color=okorange] coordinates {};
\addplot+[boxplot prepared={lower whisker=21.9, lower quartile=37.0, median=48.0, upper quartile=60.0, upper whisker=80.0}, color=okorange] coordinates {};
\addplot+[boxplot prepared={lower whisker=19.9, lower quartile=28.0, median=34.5, upper quartile=43.0, upper whisker=56.1}, color=okorange] coordinates {};
\addplot+[boxplot prepared={lower whisker=21.9, lower quartile=44.0, median=56.0, upper quartile=70.2, upper whisker=93.1}, color=okgreen] coordinates {};
\addplot+[boxplot prepared={lower whisker=22.0, lower quartile=38.0, median=50.0, upper quartile=61.2, upper whisker=81.1}, color=okblue] coordinates {};
\addplot+[boxplot prepared={lower whisker=21.0, lower quartile=40.8, median=48.0, upper quartile=61.2, upper whisker=83.4}, color=okblue] coordinates {};
\addplot+[boxplot prepared={lower whisker=20.9, lower quartile=35.8, median=45.0, upper quartile=53.0, upper whisker=78.1}, color=okblue] coordinates {};

\end{axis}
\end{tikzpicture}

\caption{Sonnet-heavy sessions made more requests for the same scripted work. Main-loop requests per session by kind of session: box from the 25th to the 75th percentile, line at
the median, whiskers from the 5th to the 95th percentile.}
\label{vt:fig:requests}
\howtoread{Boxes further right made more requests for the same scripted work. Sonnet-heavy sessions sit to the
right of the plain Opus sessions, which is part of why routing cost more on Opus.}
\end{figure}


\subsubsection{Session length}

Within the \MinTurns\ to \MaxTurns\ scripted turns of this campaign, session length moved the cost ratios little (\cref{vt:fig:band}).
On Fable the saving held at every length; on Opus the longer sessions came slightly closer to break-even.

\begin{figure}[tbp]
\centering
\pgfplotslegendfromname{bandlegend}\par\smallskip
% Cost ratio by scripted session length (all splits, 95% CIs).
\begin{tikzpicture}
\begin{groupplot}[benchmark, group style={group size=2 by 1, horizontal sep=1.6cm},
  width=0.46\linewidth, height=5.2cm, xmin=0.5, xmax=3.5, xtick={1,2,3},
  xticklabels={{$\le$8},{9--12},{$\ge$13}}, xlabel={Scripted turns}, ymajorgrids, title style={font=\small\bfseries}]
\nextgroupplot[title={Fable host}, ylabel={Cost ratio vs plain host}, ymin=0.4, ymax=0.8,
  legend to name=bandlegend, legend columns=2, legend cell align=left,
  legend style={/tikz/every even column/.append style={column sep=10pt}}]
\addplot[only marks, mark=*, mark size=2.4pt, color=okblue, error bars/.cd, y dir=both, y explicit]
  table[x=x, y=r, y error minus=m, y error plus=p]{generated/data/band-fable-sticky.dat}; \addlegendentry{sticky}
\addplot[only marks, mark=square*, mark size=2.4pt, color=okorange, error bars/.cd, y dir=both, y explicit]
  table[x=x, y=r, y error minus=m, y error plus=p]{generated/data/band-fable-shipped.dat}; \addlegendentry{shipped (re-chooses each turn)}
\nextgroupplot[title={Opus host}, ymin=0.95, ymax=1.65]
\draw[black!60] (axis cs:0.5,1) -- (axis cs:3.5,1);
\addplot[only marks, mark=*, mark size=2.4pt, color=okblue, error bars/.cd, y dir=both, y explicit]
  table[x=x, y=r, y error minus=m, y error plus=p]{generated/data/band-opus-sticky.dat};
\addplot[only marks, mark=square*, mark size=2.4pt, color=okorange, error bars/.cd, y dir=both, y explicit]
  table[x=x, y=r, y error minus=m, y error plus=p]{generated/data/band-opus-shipped.dat};
\end{groupplot}
\end{tikzpicture}

\caption{The routing ratio barely depended on scripted session length. Cost ratio by scripted session length, all splits, with 95\,\% intervals.}
\label{vt:fig:band}
\howtoread{Flat rows mean the ratio does not depend on session length within 8--16 turns. Longer sessions moved
the Opus ratios slightly toward 1; on Fable the saving held at every length.}
\end{figure}

